\begin{frame}{Tamanho e limites dependem da implementação}
\begin{itemize}
\item Um computador e um microcontrolador podem usar tamanhos diferentes para \texttt{int}.
\item \texttt{sizeof(tipo)} informa o tamanho em bytes de C.
\item \texttt{sizeof(char)} é sempre 1; a quantidade de bits nesse byte é dada por \texttt{CHAR\_BIT}.
\item \texttt{<limits.h>} fornece limites como \texttt{INT\_MIN} e \texttt{INT\_MAX}.
\item \texttt{<stdint.h>} oferece tipos de largura especificada. Os tipos exatos, como \texttt{uint16\_t}, existem quando a implementação os suporta.
\end{itemize}
\begin{block}{Decisão prática}
Consulte o ambiente de destino; não memorize um tamanho universal para \texttt{int}.
\end{block}
\end{frame}
